\chapter{Einführung}
\label{sec:einführung}

Das menschliche Fahrverhalten ist bedingt durch die Vielzahl der möglichen Situationen und Reaktionen im Straßenverkehr höchst komplex. Eine Analyse dessen ist somit nur durch initiale Vereinfachungen in einem Modell möglich oder durch eine realistische Simulation mir realen Fahrern. Für den zweiten Ansatz entwickelt das Institut für Energieeffiziente Mobilität (IEEM) ein Simulationsframework, welches mit unterschiedlichen Bewegungsmethoden (Auto, Fahrrad, Gehen) und in unterschiedlicher Realitätsnähe genutzt werden kann.

Als Softwarebasis dient das Open-Source-Projekt CARLA. CARLA ermöglicht es, verschiedene Verkehrsszenarien und Wetterbedingungen einzustellen. Die Analyse des menschlichen Verhaltens im Straßenverkehr kann somit ohne Gefahren und in verschiedensten simulierten Szenarien erfolgen. Die Server-Client Struktur von CARLA lässt ebenso mehrere Probanden in demselben Verkehrsszenario zu. So können menschliche Interaktionen realitätsnah nachgebildet werden.

Für die Analyse des Fahrverhaltens wird mindestens die Auswertung der Fahrzeugdaten (z. B.: der zurückgelegte Weg, einwirkende Kräfte) benötigt. Für eine genauere Analyse eignen sich zudem die Fahrerentscheidungen. Diese können einen Einblick geben, ob ein Fahrmanöver aus Reflex erfolgt ist oder ob es eine bewusste Entscheidung gewesen ist. Ein einfacher Weg Fahrerentscheidungen zu erfassen, ist es, dass nach einem Fahrversuch ein Fragebogen beantwortet wird. Weiterführend bietet es sich an Echtzeitdaten des Probanden aufzuzeichnen.

Eine große Rolle für das Treffen von Entscheidungen spielt die Sicht. Hierdurch wird deutlich, was die Probanden bei der Testfahrt sehen, worauf diese sich konzentrieren und welchen Fahrweg sie nehmen werden. Für die Datenerhebung der Sicht eignet sich ein sogenannter Eye-Tracker.

\section{Zielsetzung}

Im Rahmen dieser Projektarbeit soll ein kamerabasierter Eye-Tracker entwickelt werden, der in die CARLA Simulationsumgebung integriert werden kann.

Dazu werden mögliche Eye-Tracker auf die Anwendbarkeit und Integration in CARLA analysiert. Ebenfalls wird eine Simulation in CARLA erstellt, um den Eye-Tracker zu testen. Weitere Vorgaben sind die Projektion des Blickpunktes auf die Leinwand und eine Ausgabe der aufgezeichneten Daten.


\section{Lösungsansatz}

Damit die gesetzten Ziele erreicht werden, müssen zuerst die Grundlagen von Eye-Trackern und der CARLA-Simulationsumgebung erarbeitet werden. Für die Integration in das Gesamtsystem wird ein objektorientierter Python Code entwickelt. Der Code soll den Blickpunkt des Probanden als Pixel ausgeben und mit CARLA austauschen können. Ebenfalls wird die benötigte Rechenleistung minimiert. Die Ausgabe der erhobenen Daten erfolgt in Echtzeit über ein Fenster und wird im .csv-Format abgespeichert.






